\section{Canadian Engineering Competition}
\label{canadian-engineering-competition}

\subsection{Appeal Processes}
\label{appeal-processes}

\subsubsection{Appeals on Changes that Impact More Than One CEC}
\label{appeals-on-changes-that-impact-more-than-one-cec}

\begin{enumerate}
 \item
  All appeals that fit into this category must be made within two weeks of the decision by the CECAB.
 \item
  To begin the appeal process, a letter detailing the decision that is being appealed and the reason for the appeal must be submitted to the Chair of the CECAB who then must inform both the CECAB and the CFES Board.
 \item
  The CFES Board shall respond to the appeal within one week of the notice from the CECAB Chair.
 \item
  This type of appeal needs a minimum of 10 signatories on the letter from any region or at least 1 signatory from each region of the CFES.

  \begin{enumerate}
   \item
    The signatories to the letter must be the proxy holders from CFES members
  \end{enumerate}

 \item
  All decisions made by the CFES Board are final.
\end{enumerate}

\subsubsection{Appeals Impacting The Current CEC}
\label{appeals-impacting-the-current-cec}

\begin{enumerate}
 \item
  All appeals that fit into this category must be made within one week of the conclusion of CEC.
 \item
  To begin the appeal process, a letter detailing the decision that is being appealed and the reason for the appeal must be submitted to the Chair of the CECAB who then must inform both the CECAB and the CFES Board.
 \item
  The CFES Board shall respond to the appeal within one week of the notice from the CECAB Chair.
 \item
  This type of appeal needs to have consensus from the majority of the team filing the appeal, noted by the signatures of 50\%+1 of the team members on the appeal letter.

  \begin{enumerate}
   \item
    The regional head delegate may sign the letter but their signature will not count towards the percentage needed to obtain majority consensus.
  \end{enumerate}

 \item
  All decisions made by the CFES Board are final.
\end{enumerate}

\subsection{Competition Categories}
\label{competition-categories}

\subsubsection{Official Categories}
\label{official-categories}

\begin{enumerate}
 \item
  A category is considered to be an Official Category at CEC if it has been held at all Regional Qualifiers and CECs for two (2) years in a row. Once a category is considered an ``Official Category'' it must be hosted at all future CECs unless an exception is granted by the CFES Board.
 \item
  The Official Categories of CEC are:

  \begin{enumerate}
   \item
    Senior Design
   \item
    Junior Design
   \item
    Consulting Engineering
   \item
    Re-engineering
   \item
    Programming
   \item
    Innovative Design
   \item
    Extemporaneous Debate
   \item
    Engineering Communications
  \end{enumerate}

\end{enumerate}

\subsubsection{Requirements for a Category to be Hosted at CEC}
\label{requirements-for-a-category-to-be-hosted-at-cec}

\begin{enumerate}
 \item
  A category has to be hosted at all four Regional Qualifiers (WEC, OEC, QEC, AEC) or be an Official Category in order to be held at CEC.
 \item
  The decision on what categories will be hosted at CEC must be announced to the members no later than 1 week following the conclusion of the CEC of the previous year to allow members and regions time to prepare qualifiers.
\end{enumerate}
